\documentclass[aspectratio=169,10pt]{beamer}
\usepackage{fontspec}
\usepackage{polyglossia}
\setdefaultlanguage{russian}
\setotherlanguage{english}
\setmainfont{DejaVu Sans}
\setsansfont{DejaVu Sans}
\setmonofont{DejaVu Sans Mono}
\definecolor{teal}{HTML}{0A685B}
\definecolor{dark}{HTML}{153B35}
\definecolor{pale}{HTML}{E8F4EE}
\definecolor{orange}{HTML}{A95A33}
\setbeamercolor{frametitle}{fg=dark}
\setbeamercolor{title}{fg=dark}
\setbeamercolor{structure}{fg=teal}
\setbeamercolor{block title}{bg=teal,fg=white}
\setbeamercolor{block body}{bg=pale,fg=dark}
\setbeamercolor{block title alerted}{bg=orange,fg=white}
\setbeamercolor{block body alerted}{bg=orange!9,fg=dark}
\setbeamertemplate{navigation symbols}{}
\setbeamertemplate{footline}{\hfill\scriptsize\insertframenumber\, / \,\inserttotalframenumber\hspace{12pt}\vspace{5pt}}
\setbeamertemplate{itemize item}{\raisebox{.2ex}{\tiny$\blacktriangleright$}}
\setbeamerfont{frametitle}{size=\Large,series=\bfseries}
\setbeamerfont{normal text}{size=\normalsize}
\setlength{\leftmargini}{1.3em}
\begin{document}
\begin{frame}[plain]
\begin{beamercolorbox}[wd=\paperwidth,ht=\paperheight,sep=1.4cm,center]{title}
\vspace*{1.0cm}{\large Данные как основа ИИ · семестр 2}\par\vspace{1cm}
{\fontsize{30}{38}\selectfont\bfseries Очистка без потери смысла}\par\vspace{0.6cm}
{\large От диагностической таблицы к протоколу решений}\par\vfill{\small Кейс «Умный кампус» · университетские пространства}
\end{beamercolorbox}\end{frame}
\begin{frame}[t]{Сразу нельзя ответить}
\begin{itemize}
\item CO₂ = 0 ppm — прибор ничего не измерил или записал ноль?
\vspace{4pt}
\item CO₂ = 4200 ppm — ошибка выгрузки или серьёзный сигнал?
\vspace{4pt}
\item Обе строки требуют проверки; действия могут быть противоположными.
\vspace{4pt}
\end{itemize}
\vfill\begin{alertblock}{Вопрос аудитории}Какое значение удалить прямо сейчас?\end{alertblock}
\end{frame}
\begin{frame}[t]{Чему научимся}
\begin{itemize}
\item Определять единицу наблюдения и технический дубль.
\vspace{4pt}
\item Исправлять значения по доказанным правилам.
\vspace{4pt}
\item Сохранять нетипичные, но потенциально реальные события.
\vspace{4pt}
\item Проверять, как очистка повлияла на массив.
\vspace{4pt}
\end{itemize}
\end{frame}
\begin{frame}[t]{Одно наблюдение — один факт}
\begin{itemize}
\item Датчик SN-101, аудитория A-101, 16 марта, 10:00: 810 ppm.
\vspace{4pt}
\item Тот же датчик, 10:30: 810 ppm.
\vspace{4pt}
\item Одинаковые значения не означают одну и ту же физическую запись.
\vspace{4pt}
\end{itemize}
\vfill\begin{alertblock}{Вопрос аудитории}Какие поля позволяют установить тождество измерения?\end{alertblock}
\end{frame}
\begin{frame}[t]{Типы нарушений}
\begin{itemize}
\item Формат: 21,4 и 21.4; разные формы даты.
\vspace{4pt}
\item Структура: повтор ID, неизвестное помещение, конфликтные версии.
\vspace{4pt}
\item Смысл: отрицательная численность, нереалистичный нуль.
\vspace{4pt}
\item Неопределённость: экстремально высокий, но возможный CO₂.
\vspace{4pt}
\end{itemize}
\end{frame}
\begin{frame}[t]{Четыре доступных действия}
\begin{itemize}
\item Исправить по явному правилу и перепроверить.
\vspace{4pt}
\item Сохранить с пометкой качества.
\vspace{4pt}
\item Исключить конкретное значение из конкретного расчёта.
\vspace{4pt}
\item Отложить строку/поле до уточнения первоисточника.
\vspace{4pt}
\end{itemize}
\vfill\begin{alertblock}{Вопрос аудитории}Что будет, если для всех случаев выбрать удаление?\end{alertblock}
\end{frame}
\begin{frame}[t]{Пропуск не равен нулю}
\begin{itemize}
\item Нет измерения — не значит «измерено 0».
\vspace{4pt}
\item Среднее значение не восстанавливает потерянный факт.
\vspace{4pt}
\item Сохраните признак отсутствия измерения.
\vspace{4pt}
\item Проверьте, не сосредоточены ли пропуски в одном помещении.
\vspace{4pt}
\end{itemize}
\end{frame}
\begin{frame}[t]{Два вида повторов}
\begin{itemize}
\item Полная копия строки с подтверждённой повторной выгрузкой — технический дубль.
\vspace{4pt}
\item Одинаковый ID, но разные значения — конфликт, а не дубль.
\vspace{4pt}
\item Разные времена при одинаковом CO₂ — два физических наблюдения.
\vspace{4pt}
\end{itemize}
\vfill\begin{alertblock}{Вопрос аудитории}Можно ли везде использовать drop\_duplicates по CO₂?\end{alertblock}
\end{frame}
\begin{frame}[t]{Экстремум — не ошибка по определению}
\begin{itemize}
\item Случай 4200 ppm требует флага и проверки.
\vspace{4pt}
\item Если убрать высокие значения, можно скрыть аварийный режим.
\vspace{4pt}
\item Порог исключения — предметное правило, а не желание улучшить график.
\vspace{4pt}
\end{itemize}
\end{frame}
\begin{frame}[t]{Справочник помещений}
\begin{itemize}
\item A-101 и a 101 допустимо связать по документированному алиасу.
\vspace{4pt}
\item Неизвестное D-999 не присваивается ближайшей комнате.
\vspace{4pt}
\item После JOIN со справочником число наблюдений не должно незаметно вырасти или уменьшиться.
\vspace{4pt}
\end{itemize}
\vfill\begin{alertblock}{Вопрос аудитории}Почему внутреннее соединение опасно для необработанных ключей?\end{alertblock}
\end{frame}
\begin{frame}[t]{Сначала протокол, потом операция}
\begin{itemize}
\item Сохранить исходник и исходные числа строк/пропусков.
\vspace{4pt}
\item Описать правило и причину его применения.
\vspace{4pt}
\item Применить правило в копии данных.
\vspace{4pt}
\item Сверить баланс строк, ключи и изменённые поля.
\vspace{4pt}
\end{itemize}
\end{frame}
\begin{frame}[t]{ИИ предлагает. Проверка решает}
\begin{itemize}
\item ИИ может написать формулу и перечень типичных проблем.
\vspace{4pt}
\item ИИ не знает, почему датчик показал 4200 ppm.
\vspace{4pt}
\item Перед применением правил проверьте контрпример.
\vspace{4pt}
\item Никаких необратимых действий без решения человека.
\vspace{4pt}
\end{itemize}
\vfill\begin{alertblock}{Вопрос аудитории}Что в ответе ИИ является гипотезой, а не фактом?\end{alertblock}
\end{frame}
\begin{frame}[t]{Разберите плохой план ИИ}
\begin{itemize}
\item «Удалить все пропуски».
\vspace{4pt}
\item «Дубликат — любые строки с равным CO₂».
\vspace{4pt}
\item «Высокие значения CO₂ удалить как выбросы».
\vspace{4pt}
\item «Неизвестные помещения заменить ближайшими».
\vspace{4pt}
\end{itemize}
\vfill\begin{alertblock}{Вопрос аудитории}Для каждого пункта назовите риск и контроль.\end{alertblock}
\end{frame}
\begin{frame}[t]{Лабораторная: протокол исправлений}
\begin{itemize}
\item Исходные файлы: поток измерений + справочник аудиторий.
\vspace{4pt}
\item Результат: таблица, журнал решений, карантин и баланс.
\vspace{4pt}
\item Нужно показать не менее шести спорных случаев и аудит ИИ.
\vspace{4pt}
\item Работать можно в таблицах, Python или исследовательском маршруте.
\vspace{4pt}
\end{itemize}
\end{frame}
\begin{frame}[t]{Проверка качества сдачи}
\begin{itemize}
\item Единица наблюдения определена.
\vspace{4pt}
\item Каждое изменение объяснено и воспроизводимо.
\vspace{4pt}
\item Необработанные случаи видны, а не исчезли.
\vspace{4pt}
\item ИИ-предложения проверены по данным.
\vspace{4pt}
\end{itemize}
\end{frame}
\begin{frame}[t]{Выходной билет}
\begin{itemize}
\item Назовите один случай, который нельзя было исправить автоматически.
\vspace{4pt}
\item Скажите, какая проверка доказала улучшение качества данных.
\vspace{4pt}
\item Объясните, почему «пустых ячеек не осталось» не равно «данным можно доверять».
\vspace{4pt}
\end{itemize}
\end{frame}
\end{document}
